\section{Discussion and Conclusion}

We introduced \ours{}, a benchmark and sandboxed evaluation platform spanning 1,150 expert-curated tasks across eight enterprise productivity domains, with 512 tools and SQL-based verifiers authored by subject matter experts. Our experiments surface several findings that we believe have broad implications for the development of enterprise-grade LLM agents.

\textbf{Current agents are far from enterprise-ready.}
Even under oracle tool access---the most favorable possible retrieval setting---the best model achieves only 37.4\% task success, and performance degrades monotonically with horizon length and cross-domain coupling. Critically, failure modes are not random: our qualitative and quantitative analysis shows that agents struggle most with \textit{Permission and Process Compliance} (e.g., policy adherence, cascading state transitions) rather than with basic task completion. Furthermore, even frontier models refuse infeasible tasks reliably only about half the time (best: 53.9\%), falling well short of the robustness required for unsupervised deployment. These results indicate that the gap to enterprise reliability will not be closed by scaling model capacity alone.

\textbf{Planning is the dominant bottleneck, not tool execution.}
Our plan-conditioned ablations demonstrate that human-authored plans yield 14--35 percentage point gains across models and domains which is far larger than gains from automated planning or more complex multi-agent orchestration. Strikingly, Qwen3-4B conditioned on human plans is competitive with much larger models under the same condition, suggesting that once strategic reasoning is externalized, even small models can execute faithfully. This dissociation between planning and execution ability implies that the core challenge is constraint-aware plan generation, not tool invocation proficiency. Consistent with this, distractor tools do not meaningfully hurt performance further confirming that tool retrieval is not the binding constraint. Advances in long-horizon, policy-aware planning are therefore the highest-leverage direction for improving agent performance on \ours{}.

\textbf{Thinking budget matters, but has domain-specific ceilings.}
Increasing test-time compute yields substantial gains in most domains, but scaling is not universally monotonic: some domains plateau early, suggesting that additional reasoning tokens cannot compensate for fundamental gaps in domain knowledge or policy understanding. Future work should investigate how to allocate test-time compute more adaptively, and whether targeted training on constraint-heavy domains can raise these ceilings.

\textbf{Future directions.}
Our results motivate three concrete research priorities. First, \textit{constraint-aware plan generation}: methods that explicitly reason over policy constraints, side-effect dependencies, and prerequisite structures before committing to action sequences. Second, \textit{long-horizon state management}: mechanisms for maintaining coherent world state over many tool calls, such as episodic memory or structured state representations, to prevent the error accumulation we observe with increasing horizon length. Third, \textit{safe refusal and escalation}: agents must reliably detect infeasible or policy-violating requests and abstain cleanly, a capability that remains weak across all evaluated systems today.

We will release \ours{}, its sandbox environment, and evaluation tooling to support open, community-driven research. The sandbox is modular and extensible, allowing new domains, tools, and workflows to be added as enterprise practices evolve. By grounding agent evaluation in realistic, constraint-rich enterprise workflows, \ours{} aims to shift the field's focus toward the planning, safety, and policy-compliance capabilities that truly determine whether an LLM agent is deployable as a reliable \textit{AI worker}.
